\section{Support and descent for division-ring coefficients}
\label{sec:04-coefficients}

The closure construction requires two forms of algebraic coefficient
control. Disjoint subgroup support patterns must give independent
expressions. We also need descent from intersecting subgroup-field
spans when the subgroup intersection is separable in both groups.
Theorem~\ref{coeff:support} supplies these assertions.
Section~\ref{sec:analytic-localization} supplies the
separate localization argument for coefficients subject to height
translation.

Throughout this section, \(J\) is a locally indicable group and \(D_J\)
is its Hughes-free division ring over \(\Q\). For \(A\le J\), write
\(D_A\) for the division subring generated by \(\Q A\). We use the
subgroup inheritance and crossed-product properties of these fields:
if \(H\le J\) is finitely generated and
\(\chi\colon H\twoheadrightarrow\Z\) has kernel \(K\), then, choosing
\(t\in H\) with \(\chi(t)=1\),
\begin{equation}\label{coeff:cyclic-field}
 D_H=D_K(t;\operatorname{Ad}t)
 \ \hookrightarrow\ 
 D_K((t;\operatorname{Ad}t)).
\end{equation}
The left side is the Ore division ring of skew Laurent polynomials,
and the right side consists of series bounded below in degree. Distinct
subgroup cosets are independent over the corresponding subgroup field.
The normal-subgroup crossed-product assertion is
\cite[Section~2.3 and Proposition~2.2]{JaikinUniversality2021}; the
compatible rational affiliated realization is
\cite[Corollaries~6.2--6.3]{JaikinLopezAlvarez2020}, extending the
Hughes-free framework of \cite{Hughes1970}.

For arbitrary, possibly nonnormal $A$, coset independence follows
directly in that realization. If $\sum_i d_i g_i=0$ with
$d_i\in\mathcal U(A)$ and distinct right cosets $Ag_i$, put
$b=(1+\sum_i d_i d_i^*)^{-1/2}$. This is injective and all $bd_i$
are bounded. The bounded terms $bd_i g_i$ have disjoint Fourier
supports, so each is zero and hence each $d_i=0$. Left cosets use
a common right regularizer. This argument does not attribute
nonnormal coset independence to the normal-subgroup proposition.

Our arguments concern algebraic rational expressions and do not assign
Fourier coefficients to an arbitrary affiliated operator.

\subsection{Permitted support patterns}

A \emph{support pattern} is a subset of \(J\) of the form
\[
 \mathcal P=A_1g_1A_2g_2\cdots A_mg_m,
 \qquad A_i\le J,\quad g_i\in J.
\]
A product \(d_1g_1\cdots d_mg_m\), where \(d_i\in D_{A_i}\), is
\emph{permitted on} \(\mathcal P\). For \(S\subseteq J\), let
\(\mathcal V(S)\) be the \(\Q\)-span of all such products whose
support pattern is contained in \(S\). These are permitted algebraic
supports; no choice of a scalar-function realization is involved in the
definition. In particular,
\(\mathcal V(S)g=\mathcal V(Sg)\).

\begin{theorem}[Support and coset descent]\label{coeff:support}
The following assertions hold.
\begin{enumerate}[label=\textup{(\roman*)}]
\item If \(S\cap T=\varnothing\), then
\(\mathcal V(S)\cap\mathcal V(T)=0\).
\item There is a \(\Q\)-linear map
\[
 \mathcal F\colon D_J\longrightarrow\Q^J
\]
with \(\mathcal F(1)=\delta_1\), commuting with every right group
translation, such that
\(v\in\mathcal V(S)\) implies
\(\supp\mathcal F(v)\subseteq S\).
\item Let \(A,B\le J\) and \(C=A\cap B\). If \(At\cap Bs\) is empty,
then \(D_A t\cap D_B s=0\). Otherwise, for any \(x\in At\cap Bs\),
\begin{equation}\label{coeff:single-coset}
 D_A t\cap D_B s=D_Cx.
\end{equation}
If, in addition, \(C\) is separable in both \(A\) and \(B\), then
\begin{equation}\label{coeff:finite-span}
 \left(\sum_{g\in J}D_Ag\right)
 \cap
 \left(\sum_{g\in J}D_Bg\right)
 =\sum_{g\in J}D_Cg,
\end{equation}
where each sum denotes its algebraic, finite-support span.
\end{enumerate}
\end{theorem}

Applied entrywise to a finite row \(v\), the map in (ii) satisfies
\[
 \mathcal F(vR)=\mathcal F(v)R
 \qquad\text{for every compatible finite matrix }R\text{ over }\Q J.
\]
Indeed each matrix entry acts by a finite sum of right group
translations. Thus the same scalar-function map preserves the Fox
boundary and divergence equations used in the closure construction.
This consequence uses linearity and right equivariance; it requires
no multiplicativity for division-ring coefficients.

We first make the complexity argument precise. The two descent
arguments will use the same coefficient extraction, but have different
conclusions: disjoint supports give vanishing, while intersecting
single cosets give a smaller field.

\subsection{A hereditary complexity and coefficient extraction}

For predecessors of coefficient descent by rational-expression complexity,
see Dicks--Herbera--S\'anchez
\cite[Sections~3--6, especially Theorem~6.3]{DicksHerberaSanchez2004}
and Jaikin-Zapirain--L\'opez-\'Alvarez
\cite[Proposition~5.1]{JaikinLopezAlvarez2020}.
Here we combine an explicit ordinal complexity with bookkeeping for the
prescribed subgroup support patterns.

The difficulty in an induction is inversion. A coefficient of an
inverse series is a finite sum of products, but neither the lengths
of those products nor the number of their displayed terms has a
uniform bound. We therefore use two ordinal layers: the first makes
simplifying an inverse operand outweigh any finite increase in
product length, and the second makes simplifying a term outweigh
any finite increase in the number of terms. We define them on the
displayed expressions, retaining terms even when their values cancel.

A displayed rational expression can be normalized recursively as a
finite sum of terms
\begin{equation}\label{coeff:normal-term}
 q P_1^{-1}\cdots P_s^{-1}g,
 \qquad q\in\Q\setminus\{0\},\quad g\in J.
\end{equation}
Each \(P_i\) is a displayed sum in the same form, has nonzero value,
and has at least one top term with final offset \(1\). To obtain this
normalization before taking an inverse, choose a displayed top offset
\(g\), factor it to the right, and move the resulting monomial through
other inverse factors by conjugation. For example,
\(gP^{-1}=(gPg^{-1})^{-1}g\). These operations preserve the number and
nesting of inverse occurrences. They do not identify equal displayed
terms or cancel terms against one another.

The \emph{hereditary offsets} of an expression are all final group
elements appearing at any depth of its displayed expression tree.
For an inverse occurrence \(P^{-1}\), let \(H(P)\) be the subgroup
generated by the hereditary offsets of \(P\). Assign the term
\eqref{coeff:normal-term} the pattern
\begin{equation}\label{coeff:displayed-pattern}
 H(P_1)\cdots H(P_s)g.
\end{equation}
When the original factors have prescribed subgroup-field patterns,
normalize them separately before multiplying. Every resulting term
pattern is then contained in its original prescribed pattern.
Indeed all offsets used to represent an element of \(D_A\) can be
chosen in \(A\); moving a monomial through a later factor conjugates
that factor's subgroup and moves the same monomial to the right.

Write \(\oplus\) for natural, or Hessenberg, addition of ordinals:
add the coefficients of equal powers in Cantor normal form. Define the
\emph{path weight} and \emph{complexity} recursively by
\begin{align}
 w(qP_1^{-1}\cdots P_s^{-1}g)
   &=\bigoplus_{i=1}^s\omega^{c(P_i)},\label{coeff:path-weight}\\
 c(E)&=\bigoplus_{T\text{ a displayed term of }E}\omega^{w(T)}.
 \label{coeff:complexity}
\end{align}
An empty sum has complexity zero. Every complexity here is a finite
hereditary ordinal below \(\varepsilon_0\). A monomial has weight zero;
a displayed sum of \(r\) monomials has complexity \(r\). Conjugating an
expression, changing a final monomial offset, or performing the preceding
normalization preserves these complexities. When an equality is under
consideration, its joint complexity counts all displayed terms on both
sides, including terms that cancel upon evaluation.

We shall use two elementary properties of this measure. Replacing an
occurrence of \(\omega^\alpha\) in a path weight by finitely many
\(\omega^{\beta_i}\), with every \(\beta_i<\alpha\), strictly decreases
that path weight. Replacing a displayed term of weight \(\eta\) by
finitely many terms of weights less than \(\eta\) strictly decreases
the outer complexity. Both statements follow directly from Cantor
normal form; unaffected summands may be retained by strict monotonicity
of natural addition. There is no bound on the finite number of replacement
terms. Also, the complexity of any inverse operand in a displayed
expression is strictly smaller than the complexity of the expression
containing it.

Relative to a character \(\chi\), call an inverse occurrence
\(P^{-1}\), or its normalized operand \(P\), \emph{nonhomogeneous}
when \(\chi\) is nonzero on some hereditary offset of \(P\),
equivalently on \(H(P)\). This refers to the displayed expression:
for example, \(1+t-t\) is nonhomogeneous when \(\chi(t)\ne0\),
although its value is constant. When all hereditary offsets are
killed by \(\chi\), the operand can instead be retained unchanged
in degree zero.

\begin{lemma}[Coefficient extraction invariant]\label{coeff:extraction}
Let finitely many normalized rational expressions have hereditary
offsets in a finitely generated subgroup \(H\), and let
\(\chi\colon H\twoheadrightarrow\Z\) be a character with kernel
\(K\). Expand in the positive skew Laurent field in
\eqref{coeff:cyclic-field}, writing powers of \(t\) on the right.
Every coefficient has a finite normalized rational expression in
\(D_K\), with the following properties.
\begin{enumerate}[label=\textup{(\alph*)}]
\item If a summand is permitted on \(S\), the expression for its
coefficient of degree \(n\) has term patterns contained in
\[
 (S\cap\chi^{-1}(n))t^{-n}.
\]
Here \(S\) can be replaced by its intersection with \(H\).
\item Coefficient extraction does not increase the complexity of an
arbitrary top expression. If an inverse occurrence is nonhomogeneous,
the decrease is strict. In particular, for a normalized operand \(P\)
whose hereditary offsets are not all killed by \(\chi\), every
coefficient expression has complexity strictly less than \(c(P)\).
\item For a collection of top expressions, assume at least one
represented top offset equals \(1\) and that \(\chi\) is nonzero on
the subgroup generated by all their hereditary offsets. The joint
complexity of the coefficient expressions at each degree is strictly
smaller than the original joint complexity.
\item If the input is represented in one field coset \(D_Aa\), with
all inverse operands internal to \(A\), a nonzero coefficient of
degree \(n\) is represented in one coset of \(D_{A\cap K}\). Its
permitted group coset is
\((Aa\cap\chi^{-1}(n))t^{-n}\), and its inverse operands remain
internal to \(A\cap K\).
\end{enumerate}
\end{lemma}

\begin{proof}
We construct the coefficients by induction on the expression trees.
An inverse operand whose hereditary offsets are all killed by
\(\chi\) is retained unchanged in degree zero. For any other inverse
operand \(P\), first use a transporter inside \(H(P)\): if
\(\chi(H(P))=d\Z\), choose \(a\in H(P)\) with \(\chi(a)=d>0\).
The cyclic-field property for this subgroup gives the expansion of
\(P\) in powers of \(a\). Its leading nonzero coefficient is itself
a rational expression. Factoring this coefficient and the leading
power gives the inverse expansion by the geometric series
\((1+R)^{-1}=\sum_{j\ge0}(-R)^j\), where \(R\) has positive
valuation. A fixed coefficient of that expansion is a finite sum of
finite products of coefficients of \(P\) and inverses of its leading
coefficient, with the conjugations imposed by the skew multiplication.
All series in this construction are bounded below, so coefficient
multiplication and inverse expansion involve only finitely many terms
in each fixed degree.

For completeness, the complexity comparison can be proved at the same
time by structural induction. A normalized nonhomogeneous operand
\(P\) has a top offset \(1\). If one of its inverse occurrences is
nonhomogeneous, the induction hypothesis decreases the complexities of
that occurrence's operand coefficients. If all its inverse occurrences
are homogeneous, every term contributes to just one degree, and the
nonzero character distinguishes at least two top offsets. Thus every
degree omits at least one term. These are the two cases proving that
every coefficient of \(P\) has complexity less than \(c(P)\).
Consequently, in expanding \(P^{-1}\), the inverse of its leading
coefficient has operand complexity less than \(c(P)\). Every inverse
operand appearing in one of the other coefficients also has complexity
less than \(c(P)\). Any coefficient term of \(P^{-1}\) therefore has
path weight less than \(\omega^{c(P)}\). This replaces the original
inverse's entire path contribution by finitely many lower contributions.
The two natural-sum comparisons above prove strict decrease even after
arbitrary finite multiplication and distribution of the coefficient
expressions. They also justify the nonhomogeneous-inverse case used
in the preceding structural induction.

This proves (b). For (c), retain the tags recording each original term.
If some inverse occurrence is nonhomogeneous, every resulting term
coming from that occurrence has smaller path weight, while unchanged
terms contribute at most once at any degree. If all inverse occurrences
are homogeneous, the top offset \(1\), together with nontriviality on
all hereditary offsets, forces another top degree. Every coefficient
then omits a top term. This proves (c), also for a joint collection
whose individual expressions need not have a top offset \(1\).

It remains to check that this construction respects the designated
supports rather than just the field identities. In the local expansion
at \(P^{-1}\), all transporters lie in \(H(P)\), and its coefficient
of height \(n\) is supported on the corresponding coset of
\(H(P)\cap K\). To compare with the global variable, write
\(a^m=k_m t^{dm}\), with \(k_m\in K\). A local coefficient then
acquires the final monomial offset \(k_m\). Moving powers of \(t\)
to the right in a product conjugates its subsequent kernel factors.
Multiplying the corresponding group elements shows that every resulting
term pattern is contained in the original product pattern's total-degree
slice, followed by \(t^{-n}\). The new offsets are monomials; their word
lengths play no role in \eqref{coeff:path-weight}--\eqref{coeff:complexity}.
Conjugations and final-offset changes preserve the measure. This proves
(a) and its compatibility with (b) and (c). For (d), all the local
transporters before the final offset lie in \(A\). When intermediate
global powers are moved to the right and the resulting monomials are
normalized, each conjugating prefix is consequently an element of
\(A\). It preserves \(A\cap K\), and the final offset lies in the
displayed single coset. This also follows by first expanding the whole
\(D_A\)-factor using an \(A\)-transporter and only then converting its
coefficients to the global variable.
\end{proof}

\subsection{Disjoint supports and the scalar-function map}

\begin{proof}[Proof of Theorem~\ref{coeff:support}(i)--(ii)]
Suppose \(v\in\mathcal V(S)\cap\mathcal V(T)\), with \(S\cap T\)
empty. Choose finite rational representations of its two values,
retaining the prescribed pattern for each top summand. We induct on
their joint complexity. Right-multiply the equality by the inverse of
one represented top offset. This translates both allowed sets and
preserves their disjointness and the joint complexity, while putting
one top offset at \(1\). Let \(H\) be the finitely generated subgroup
of all hereditary offsets in the shifted equality.

If \(H=1\), all nonzero displayed terms have pattern \(\{1\}\).
Disjointness forces one side to have no such term, and the equality
therefore has value zero. Otherwise local indicability gives an
epimorphism \(\chi\colon H\twoheadrightarrow\Z\). Apply
Lemma~\ref{coeff:extraction}. At each degree the two coefficient
expressions have strictly smaller joint complexity, and their allowed
supports are disjoint subsets of \(K\). The induction hypothesis
makes each side's coefficient zero. Injectivity of
\eqref{coeff:cyclic-field} makes both original values zero as well.
Well-foundedness of the ordinal measure completes the induction.

Set \(W=\mathcal V(J\setminus\{1\})\). Part (i) gives
\(W\cap\Q1=0\). Define a linear functional on \(W+\Q1\) by
\(\ell(w+q)=q\), and extend it to a \(\Q\)-linear functional on
\(D_J\). Put
\begin{equation}\label{coeff:scalar-map}
 \mathcal F(v)(g)=\ell(vg^{-1}).
\end{equation}
Then \(\mathcal F(1)=\delta_1\). If \(g\notin S\) and
\(v\in\mathcal V(S)\), then
\(vg^{-1}\in\mathcal V(Sg^{-1})\subseteq W\), so this coefficient
vanishes. Finally,
\[
 \mathcal F(vh)(g)=\ell(vhg^{-1})
   =\mathcal F(v)(gh^{-1}),
\]
which is precisely equivariance for right group multiplication.
This proves all support restrictions simultaneously. No multiplicativity
of \(\mathcal F\) is asserted or needed.
\end{proof}

\subsection{Rational descent from a smaller coefficient field}

We isolate the rational-series step, including its handedness.

\begin{lemma}[Rational descent]\label{coeff:rational-descent}
Let \(E\subseteq F\) be division rings and let conjugation by \(u\)
induce an automorphism of \(F\) preserving \(E\). If an element of
\(F(u)\) has a bounded-below Laurent expansion
\[
 S=\sum_n u^n f_n,\qquad f_n\in E,
\]
then \(S\in E(u)\).
\end{lemma}

\begin{proof}
Choose a nonzero right Laurent-polynomial denominator
\(P=\sum_{j\in I}u^jp_j\), with \(I\) finite and \(p_j\in F\),
such that \(SP\) is a Laurent polynomial. For all sufficiently large
\(n\), its coefficient of \(u^n\), with coefficients written on the
right, is zero:
\begin{equation}\label{coeff:recurrence}
 \sum_{j\in I}(u^{-j}f_{n-j}u^j)p_j=0.
\end{equation}
Every parenthesized coefficient belongs to \(E\). Choose a basis
\((b_\gamma)_\gamma\) of \(F\) as a \emph{left} vector space over
\(E\), and expand the finitely many coefficients as
\(p_j=\sum_\gamma e_{j\gamma}b_\gamma\). Only finitely many basis
vectors are used. Linear independence makes every coordinate tuple
\((e_{j\gamma})_{j\in I}\) satisfy all equations
\eqref{coeff:recurrence} simultaneously. Since \(P\ne0\), at least
one such tuple is nonzero. Fix it and put
\(P'=\sum_{j\in I}u^je_{j\gamma}\ne0\).

The product \(SP'\) has coefficients in \(E\), has a lower degree
bound, and vanishes in every sufficiently high degree. It is therefore
a Laurent polynomial over \(E\), and \(S=(SP')(P')^{-1}\in E(u)\).
There is no finite-subsystem argument here: one fixed coordinate tuple
satisfies the entire family of recurrence equations.
\end{proof}

\subsection{Intersections of field cosets}

\begin{proof}[Proof of Theorem~\ref{coeff:support}(iii), single cosets]
The empty-intersection assertion follows from part (i). Suppose
\(At\cap Bs=Cx\) and \(0\ne v\in D_At\cap D_Bs\). We induct on the
joint complexity of expressions for \(v\) in these two field cosets.
Choose all inverse operands on the first side inside \(A\), and all
those on the second inside \(B\). Right-shift by one represented top
offset as in the preceding proof, and let \(H\) be the subgroup
generated by the new hereditary offsets.

Intersect the two permitted cosets with \(H\). Each nonempty
intersection is a coset of \(A\cap H\), respectively \(B\cap H\).
All represented inverse operands remain internal to these subgroups:
their offsets lie in the original subgroup and, by definition, in
\(H\). Likewise, ratios of represented top offsets on a fixed side
lie in its intersected subgroup. Thus this restriction neither changes
the expression nor loses its single-coset form. The restricted cosets
have a common point by part (i), since their common field value is
nonzero. We may accordingly work entirely in \(H\), temporarily
reusing \(A,B,C\) for their intersections with \(H\).
If \(H=1\), the assertion is immediate.
Otherwise choose \(\chi\colon H\twoheadrightarrow\Z\), with kernel
\(K\).

For each degree \(n\), the two coefficient expressions lie in single
cosets of the fields of \(A\cap K\) and \(B\cap K\), allowing the
final offsets in \(K\) described in
Lemma~\ref{coeff:extraction}. Their joint complexity has strictly
decreased. The induction hypothesis therefore says that their common
value is zero if the two group cosets do not meet, and otherwise lies
in the field of \(A\cap B\cap K\) times a point of their
intersection. Returning to the original degree, the coefficient is
thus permitted precisely on the corresponding slice of the common
coset \(Cx\), with its subgroup-field coefficient in
\(D_{C\cap K}\).

Choose a common point \(x\) inside \(H\) and divide \(v\) by it on
the right. The resulting rational series has coefficients in the common
kernel field and group support restricted to \(C\). If
\(\chi(C)=0\), it has only degree zero and already belongs to
\(D_{C\cap K}=D_C\). If \(\chi(C)=d\Z\), with \(d>0\), choose
\(u\in C\) of height \(d\), and put
\(E=D_{C\cap K}\), \(F=D_K\). Then it has the form
\begin{equation}\label{coeff:common-series}
 vx^{-1}=\sum_m u^m f_m,\qquad f_m\in E,
\end{equation}
and conjugation by \(u\) preserves \(E\).

We must also check that \eqref{coeff:common-series} is rational over
\(F(u)\), not just over \(F(t)\). The subgroup
\(H_d=\chi^{-1}(d\Z)\) has index \(d\) in \(H\), and its field is
\(F(u)\). The finite crossed decomposition gives
\[
 D_H=\bigoplus_{r=0}^{d-1}D_{H_d}t^r.
\]
The Laurent series of the \(r\)-th summand uses only degrees congruent
to \(r\) modulo \(d\). Since \eqref{coeff:common-series} has only
zero-congruence degrees, injectivity of Laurent expansion makes all
other summands vanish. Hence \(vx^{-1}\in F(u)\).
Lemma~\ref{coeff:rational-descent} now gives
\(vx^{-1}\in E(u)=D_C\). This proves the required inclusion; the
reverse one follows from \(C\le A,B\). Undoing the initial common
right shift proves \eqref{coeff:single-coset} in the original frame.
If a restricted ambient subgroup was used, its common subgroup is
contained in the original \(C\), so this conclusion also gives the
original claim.
\end{proof}

\subsection{Finite spans and separability}

\begin{proof}[Proof of Theorem~\ref{coeff:support}(iii), finite spans]
One inclusion in \eqref{coeff:finite-span} is immediate. For the
other, write a common element as
\[
 v=\sum_{i=1}^{r}a_i t_i=\sum_{j=1}^{s}b_j s_j,
 \qquad a_i\in D_A,\quad b_j\in D_B.
\]
Merge repeated cosets on each side, so the original \(At_i\)'s are
pairwise disjoint, as are the \(Bs_j\)'s. Every nonempty overlap is
\(At_i\cap Bs_j=Cx_{ij}\).

For a fixed \(i\), two distinct overlaps \(Cx_{ij}\) and
\(Cx_{ik}\) have
\(x_{ij}x_{ik}^{-1}\in A\setminus C\). Separability supplies a
finite-index subgroup \(A_0\le A\) containing \(C\) and excluding
all these finitely many differences. Hence no right \(A_0\)-coset
inside \(At_i\) meets two distinct original \(B\)-cosets. To see
this independently of representatives, a common \(A_0\)-coset meeting
both overlaps would imply
\(c x_{ij}(c'x_{ik})^{-1}\in A_0\) for some \(c,c'\in C\), and
therefore \(x_{ij}x_{ik}^{-1}\in A_0\), a contradiction. Choose
\(B_0\le B\) symmetrically, of finite index and containing \(C\).

These two choices already give the simultaneous refinement required.
Since \(A_0\cap B=C\), the intersection of a refined \(A_0\)-coset
with its unique possible original \(B\)-coset is one \(C\)-coset.
As \(C\le B_0\), that entire intersection lies in just one refined
\(B_0\)-coset. The symmetric argument applies on the other side.
Thus the bipartite intersection graph of all refined cosets is a
matching, with possible isolated vertices. Also
\(A_0\cap B_0=C\).

We use the finite-index field decompositions
\[
 D_A=\bigoplus_{a\in A_0\backslash A}D_{A_0}a,
 \qquad
 D_B=\bigoplus_{b\in B_0\backslash B}D_{B_0}b.
\]
For clarity, these do not require the chosen subgroups to be normal.
One may first pass to a finite-index normal core: its finite crossed
span in the overfield is a finite-dimensional domain over its division
ring, hence a division ring, and so equals the larger subgroup field.
Grouping the resulting summands by the original subgroup cosets gives
the displayed formulas.

Refine the two expressions for \(v\) by these decompositions. Group
the terms of their difference according to the connected components
of the refined intersection graph. Distinct components have disjoint
unions of permitted supports. Part (i) therefore makes the value of
each component zero. Isolated terms vanish. Each remaining component
is a single pair, giving an equality between one \(D_{A_0}\)-coset
and one \(D_{B_0}\)-coset. By \eqref{coeff:single-coset}, their common
value belongs to \(D_Cx\) for a point of their group-coset overlap.
Adding the finitely many pairs proves \eqref{coeff:finite-span}.
\end{proof}
